\chapter{Binary Search Trees}
\chapterepigraph{A binary search tree is binary search wearing a tree's
clothing. Every node carries the same promise an array's sortedness
carries -- everything in one direction is smaller, everything in the other
is larger -- and that single promise is enough to make search, insertion,
ordered traversal, and rank queries all run in time proportional to the
tree's height, not its size.}

\section{What Makes a Tree a Search Tree?}
\label{sec:bst-intro}

A \textbf{binary search tree (BST)} is a binary tree satisfying the
\textbf{BST invariant}: for every node, every value in its left subtree
is strictly less than the node's own value, and every value in its right
subtree is strictly greater. This is a direct generalization of the
sorted-array invariant that powered the entire Binary Search chapter,
restructured so that ``go left'' or ``go right'' replaces ``halve the
search window.''

\begin{ideabox}[The One Fact That Powers Almost Every Section in This Chapter]
An \textbf{inorder traversal} (left subtree, then the node, then right
subtree) of a BST always visits nodes in \textbf{strictly increasing
order}. This is not a coincidence -- it is a direct restatement of the BST
invariant applied recursively at every node, and it is the single fact
that connects BSTs back to everything already proven about sorted arrays
in earlier chapters: an inorder traversal of a BST \emph{is}, in effect, a
sorted array, just discovered one $O(\log n)$-depth recursive step at a
time instead of stored contiguously.
\end{ideabox}

\subsection*{How to Recognise Which Technique a BST Problem Needs}

\begin{patternbox}[Recurring BST Keywords]
\begin{itemize}
  \item \textbf{``Search/insert/delete a value''} -- a single
        root-to-leaf walk, choosing left or right at each node by
        comparing to the BST invariant, exactly like binary search
        choosing which half to discard.
  \item \textbf{``Kth smallest/largest,'' ``floor/ceil,'' or any other
        \emph{rank}-based query''} -- either an inorder traversal (which
        visits values in sorted order for free) or a single guided
        root-to-leaf walk that narrows toward the answer using the BST
        invariant directly, without visiting every node.
  \item \textbf{``Validate,'' ``two nodes are swapped,'' or any property
        phrased in terms of \emph{sortedness}''} -- think in terms of the
        inorder sequence being (or not being) sorted, rather than
        reasoning node-by-node about left/right children directly.
  \item \textbf{``Predecessor/successor,'' or ``iterator that returns
        values one at a time in sorted order''} -- a controlled,
        space-efficient version of inorder traversal, often using an
        explicit stack to simulate the recursion only as deep as
        currently needed.
  \item \textbf{``Largest BST within a general binary tree''} -- a
        bottom-up recursive technique returning, from each subtree,
        enough information (is it a valid BST? what is its size? what
        are its min and max?) for its parent to decide validity and size
        in O(1) extra work per node, very much in the spirit of dynamic
        programming on trees.
\end{itemize}
\end{patternbox}

Each section in this chapter follows the established structure: problem
statement and keywords; the reasoning that identifies the right traversal
or recursive technique and why it is correct; the algorithm in words and
pseudocode; a hand-traced dry run; complete compilable C++ code; and a
time/space complexity analysis.

\newpage

\section{Search in a Binary Search Tree}
\label{sec:bst-search}

\problemstatement{Given the root of a binary search tree and a target
value $X$, find and return the subtree rooted at the node whose value
equals $X$, or return null if no such node exists.}

\begin{patternbox}
\emph{``Search for a value in a BST''} is the tree-shaped restatement of
ordinary binary search (the Binary Search chapter's opening section): the BST invariant
plays exactly the role sortedness played there, letting a single
comparison at each node discard an entire subtree (half the remaining
candidates) rather than requiring us to check it at all.
\end{patternbox}

\subsection*{Understanding the problem carefully}

At any node, comparing $X$ to that node's value tells us everything we
need: if $X$ equals the node's value, we are done. If $X$ is smaller, the
BST invariant guarantees that \emph{every} value in the right subtree is
even larger than the current node, so $X$ cannot possibly be there --
discard the right subtree entirely and recurse left. Symmetrically, if
$X$ is larger, discard the left subtree and recurse right.

\begin{ideabox}[One Comparison, One Subtree Discarded]
At each node: if the node is null, $X$ is not present (return null). If
$X$ equals the node's value, return this node. If $X<$ the node's value,
recurse into the left child. If $X>$ the node's value, recurse into the
right child.
\end{ideabox}

\subsection*{Why this is exactly binary search, with depth playing the role of $\log n$ steps}

Each step travels one level deeper into the tree and provably eliminates
an entire subtree from consideration, using the same discard argument as
array-based binary search: the eliminated subtree's values are
\emph{all} on the wrong side of $X$, guaranteed by the BST invariant, so
none of them need to be checked individually. Because each step descends
exactly one level, the total number of steps is bounded by the tree's
height -- $O(\log n)$ for a balanced tree, though note this degrades to
$O(n)$ for a degenerate (linked-list-shaped) tree, a distinction worth
remembering throughout this chapter.

\subsection*{Algorithm}

\begin{enumerate}
  \item If the current node is null: return null.
  \item If $X =$ current node's value: return the current node.
  \item If $X <$ current node's value: recurse into the left child.
  \item Else: recurse into the right child.
\end{enumerate}

\begin{algorithm}[H]
\caption{Search in a BST}
\begin{algorithmic}[1]
\Procedure{SearchBST}{\textit{root}, $X$}
  \If{$\textit{root} = \text{null}$} \Return \text{null} \EndIf
  \If{$X = \textit{root}.\textit{val}$} \Return \textit{root} \EndIf
  \If{$X < \textit{root}.\textit{val}$}
    \State \Return \Call{SearchBST}{$\textit{root}.\textit{left}$, $X$}
  \Else
    \State \Return \Call{SearchBST}{$\textit{root}.\textit{right}$, $X$}
  \EndIf
\EndProcedure
\end{algorithmic}
\end{algorithm}

\subsection*{Dry run}

\dryrun{Take the BST
$\begin{smallmatrix} & & 5 & & \\ & 3 & & 8 & \\ 1 & & 4 & & 9
\end{smallmatrix}$ (root $5$; left child $3$ with children $1,4$; right
child $8$ with right child $9$), searching for $X=4$.}

\begin{itemize}
  \item At $5$: $4<5$, go left to $3$.
  \item At $3$: $4>3$, go right to $4$.
  \item At $4$: $4=4$, found. Return this node.
\end{itemize}

\subsection*{C++ implementation}

\begin{lstlisting}
#include <bits/stdc++.h>
using namespace std;

struct TreeNode {
    int val;
    TreeNode* left;
    TreeNode* right;
    TreeNode(int v) : val(v), left(nullptr), right(nullptr) {}
};

TreeNode* searchBST(TreeNode* root, int X) {
    if (root == nullptr || root->val == X) return root;
    if (X < root->val) return searchBST(root->left, X);
    return searchBST(root->right, X);
}

int main() {
    TreeNode* root = new TreeNode(5);
    root->left = new TreeNode(3);
    root->right = new TreeNode(8);
    root->left->left = new TreeNode(1);
    root->left->right = new TreeNode(4);
    root->right->right = new TreeNode(9);

    TreeNode* found = searchBST(root, 4);
    cout << (found ? found->val : -1) << endl;
    return 0;
}
\end{lstlisting}

\begin{complexitybox}
\textbf{Time Complexity.} $O(h)$, where $h$ is the tree's height --
$O(\log n)$ for a balanced BST, $O(n)$ in the worst case for a
degenerate, linked-list-shaped tree.
\textbf{Space Complexity.} $O(h)$ for the recursion stack (or $O(1)$ if
rewritten iteratively, since no backtracking is ever needed).
\end{complexitybox}

\begin{takeawaybox}
Every BST operation in this chapter that walks from the root toward a
single target (search, insert, floor/ceil, predecessor/successor)
inherits this exact $O(h)$ complexity and this exact ``one comparison,
one subtree discarded'' justification. Internalizing it once here saves
re-deriving it in every later section.
\end{takeawaybox}

\section{Find the Minimum and Maximum in a BST}
\label{sec:bst-min-max}

\problemstatement{Given the root of a BST, find its minimum value and its
maximum value.}

\begin{patternbox}
\emph{``Minimum/maximum of a BST''} requires no comparison logic at all,
unlike Section~\ref{sec:bst-search} -- the BST invariant alone pins the
minimum and maximum to two fixed, unambiguous locations: the
\textbf{leftmost} node and the \textbf{rightmost} node.
\end{patternbox}

\subsection*{Understanding the problem carefully}

By the BST invariant, every left child is smaller than its parent, and
every right child is larger. So starting from the root and always moving
left strictly decreases the value at every step, until no left child
remains -- at which point we are standing on a value smaller than every
other value in the tree (every node we passed is larger, and every node in
their right subtrees, which we skipped, is even larger still). The
symmetric argument, always moving right, finds the maximum.

\begin{ideabox}[Walk to the Edge, No Comparisons Needed]
For the minimum: starting at the root, repeatedly move to the left child
until there is none; return the value at that point. For the maximum:
symmetrically, repeatedly move right until there is none.
\end{ideabox}

\subsection*{Why no value comparisons are needed during the walk}

This is a direct structural consequence of the BST invariant, not
something that needs separate verification at each step: the invariant
\emph{guarantees}, for every node visited, that its left child (if any) is
strictly smaller, so simply choosing ``left, if it exists'' at every step
is already provably optimal -- there is no need to compare the current
node's value to anything, since the tree's shape alone already encodes
the answer's location.

\subsection*{Algorithm}

\begin{enumerate}
  \item \textsc{FindMin}: starting at \textit{root}, while the current
        node has a left child, move to it. Return the current node's
        value.
  \item \textsc{FindMax}: symmetrically, while the current node has a
        right child, move to it. Return the current node's value.
\end{enumerate}

\subsection*{Dry run}

\dryrun{Take the BST with root $5$, left child $3$ (with left child $1$
and right child $4$), right child $8$ (with right child $9$).}

\textbf{Min}: $5 \to 3 \to 1$ (no left child). Minimum is $1$.

\textbf{Max}: $5 \to 8 \to 9$ (no right child). Maximum is $9$.

\subsection*{C++ implementation}

\begin{lstlisting}
#include <bits/stdc++.h>
using namespace std;

struct TreeNode {
    int val;
    TreeNode* left;
    TreeNode* right;
    TreeNode(int v) : val(v), left(nullptr), right(nullptr) {}
};

int findMin(TreeNode* root) {
    while (root->left) root = root->left;
    return root->val;
}

int findMax(TreeNode* root) {
    while (root->right) root = root->right;
    return root->val;
}

int main() {
    TreeNode* root = new TreeNode(5);
    root->left = new TreeNode(3);
    root->right = new TreeNode(8);
    root->left->left = new TreeNode(1);
    root->left->right = new TreeNode(4);
    root->right->right = new TreeNode(9);

    cout << "Min: " << findMin(root) << ", Max: " << findMax(root) << endl;
    return 0;
}
\end{lstlisting}

\begin{complexitybox}
\textbf{Time Complexity.} $O(h)$, where $h$ is the tree's height.
\textbf{Space Complexity.} $O(1)$ -- the iterative walk needs no
recursion stack at all, unlike most other BST operations in this chapter.
\end{complexitybox}

\begin{takeawaybox}
Finding the minimum and maximum of a BST is the simplest possible
operation in this chapter precisely because it needs zero comparisons --
the structure alone answers the question. This ``walk to the leftmost /
rightmost node'' primitive is reused directly inside Delete
(Section~\ref{sec:bst-delete}) to find a replacement node, and inside
Predecessor/Successor (Section~\ref{sec:bst-pred-succ}).
\end{takeawaybox}

\section{Insert a Node in a BST}
\label{sec:bst-insert}

\problemstatement{Given the root of a BST and a value $X$ to insert,
insert $X$ into the tree so that the BST invariant still holds, and
return the (possibly new) root.}

\begin{patternbox}
\emph{``Insert into a BST''} reuses Section~\ref{sec:bst-search}'s
search walk almost unchanged: insertion is what happens when a search for
$X$ runs off the edge of the tree (reaches a null pointer) instead of
finding a match -- that null position is exactly where $X$ belongs.
\end{patternbox}

\subsection*{Understanding the problem carefully}

Walk down the tree exactly as in Section~\ref{sec:bst-search}'s search:
at each node, go left if $X$ is smaller, right if $X$ is larger. Since
$X$ is not yet in the tree, this walk will never find an exact match --
it will instead reach a null child pointer. That null position is the
unique place where $X$ can be attached while preserving the BST
invariant: it is smaller than everything on the path so far that went
right, and larger than everything that went left, by construction of the
walk itself.

\begin{ideabox}[Same Walk as Search, Attach at the Null]
Walk from the root as in a search for $X$. Whenever the walk would move
into a null child, create a new node with value $X$ and attach it there
instead. If the tree was empty to begin with, the new node simply becomes
the root.
\end{ideabox}

\subsection*{Why attaching at the first null reached is always valid}

Every step of the walk maintains an invariant: $X$ is consistent with
becoming a descendant of the current node, given everything decided so
far (it is greater than every ancestor we went right from, and less than
every ancestor we went left from). Reaching a null pointer means no
existing node currently occupies the position consistent with all those
decisions -- so creating a new node there extends the tree without
violating the BST invariant anywhere, since the new node's relationship
to every ancestor on the path is already exactly correct by construction.

\subsection*{Algorithm}

\begin{enumerate}
  \item If the current node is null: create and return a new node with
        value $X$.
  \item If $X <$ current node's value: recursively insert into the left
        child, and attach the (possibly updated) result back as the left
        child.
  \item Else: recursively insert into the right child, attach back as
        the right child.
  \item Return the current node (unchanged, except possibly for the
        newly attached child).
\end{enumerate}

\begin{algorithm}[H]
\caption{Insert into a BST}
\begin{algorithmic}[1]
\Procedure{InsertBST}{\textit{root}, $X$}
  \If{$\textit{root} = \text{null}$}
    \State \Return new node with value $X$
  \EndIf
  \If{$X < \textit{root}.\textit{val}$}
    \State $\textit{root}.\textit{left} \gets \Call{InsertBST}{\textit{root}.\textit{left}, X}$
  \Else
    \State $\textit{root}.\textit{right} \gets \Call{InsertBST}{\textit{root}.\textit{right}, X}$
  \EndIf
  \State \Return \textit{root}
\EndProcedure
\end{algorithmic}
\end{algorithm}

\subsection*{Dry run}

\dryrun{Insert $X=6$ into the BST with root $5$, left child $3$, right
child $8$.}

\begin{itemize}
  \item At $5$: $6>5$, go right to $8$.
  \item At $8$: $6<8$, go left -- but $8$ has no left child (null).
        Create new node $6$, attach as $8$'s left child.
\end{itemize}

The tree now has $8$'s left child equal to $6$.

\subsection*{C++ implementation}

\begin{lstlisting}
#include <bits/stdc++.h>
using namespace std;

struct TreeNode {
    int val;
    TreeNode* left;
    TreeNode* right;
    TreeNode(int v) : val(v), left(nullptr), right(nullptr) {}
};

TreeNode* insertBST(TreeNode* root, int X) {
    if (root == nullptr) return new TreeNode(X);

    if (X < root->val) root->left = insertBST(root->left, X);
    else root->right = insertBST(root->right, X);

    return root;
}

int main() {
    TreeNode* root = new TreeNode(5);
    root->left = new TreeNode(3);
    root->right = new TreeNode(8);

    insertBST(root, 6);
    cout << "8's left child: " << root->right->left->val << endl;
    return 0;
}
\end{lstlisting}

\begin{complexitybox}
\textbf{Time Complexity.} $O(h)$, identical to search.
\textbf{Space Complexity.} $O(h)$ for the recursion stack, plus $O(1)$
for the single new node created.
\end{complexitybox}

\begin{takeawaybox}
Insertion is search that creates a node instead of reporting failure --
recognizing this means you do not need to separately re-derive the walk
logic, only add the ``attach a new node here'' step at the point where
search would otherwise give up.
\end{takeawaybox}

\section{Delete a Node in a BST}
\label{sec:bst-delete}

\problemstatement{Given the root of a BST and a value $X$, delete the
node with value $X$ (if present) and return the root of the resulting
BST, which must still satisfy the BST invariant.}

\begin{patternbox}
\emph{``Delete from a BST''} is the hardest of the three basic operations
because, unlike search or insert, the node being removed might be in the
\textbf{middle} of the tree with two children that both need to stay
connected to the rest of the tree somehow. The keyword to notice: deletion
needs a \textbf{replacement value} that preserves the BST invariant, and
Section~\ref{sec:bst-min-max}'s ``walk to the edge'' primitive is exactly
what finds one.
\end{patternbox}

\subsection*{Understanding the problem carefully}

First, locate the node with value $X$ using the same walk as
Section~\ref{sec:bst-search}. Once found, there are three structurally
different cases:

\begin{itemize}
  \item \textbf{Leaf node (no children).} Simply remove it -- nothing
        else needs to change.
  \item \textbf{One child.} Replace the node with its single child --
        the parent's pointer to this node is redirected to point at the
        child instead, and the BST invariant is automatically preserved,
        since that child's entire subtree already sat correctly on one
        side of the now-removed node.
  \item \textbf{Two children.} This is the hard case: we cannot just
        remove the node, since that would disconnect two subtrees from
        each other. Instead, find a \textbf{replacement value} that can
        legally take this node's place -- the smallest value in the
        right subtree (its inorder successor) or, symmetrically, the
        largest value in the left subtree (its inorder predecessor) both
        work. Copy that replacement value into the current node, then
        recursively delete the replacement's \emph{original} location
        from whichever subtree it came from (which, being an edge value
        of that subtree, will only ever fall into the leaf or one-child
        case, never back into the two-children case).
\end{itemize}

\begin{ideabox}[Two-Children Case: Borrow the Successor, Then Delete It Recursively]
When deleting a node with two children, find its inorder successor (the
minimum of its right subtree, found via Section~\ref{sec:bst-min-max}'s
walk), copy that successor's value into the current node, and then
recursively delete that same value from the right subtree -- which now
contains a duplicate, but only temporarily, and the recursive delete call
removes the original occurrence cleanly.
\end{ideabox}

\subsection*{Why borrowing the successor's value, then deleting it, is always correct}

The inorder successor of a two-children node is, by definition, the
smallest value strictly greater than the node's current value -- which
means it is simultaneously larger than everything in the node's left
subtree (since it's larger than the node itself) and smaller than or
equal to everything else remaining in the right subtree (since it's the
minimum of that subtree). So placing it directly into the current node's
position preserves the BST invariant relative to both subtrees perfectly.
And because the successor is the minimum of the right subtree, it has no
left child (a node with a left child cannot be a subtree's minimum), so
deleting its original location is guaranteed to fall into the simple
leaf-or-one-child case, never recursing into another two-children
scenario at that step.

\subsection*{Algorithm}

\begin{enumerate}
  \item If the current node is null: return null ($X$ was not found;
        nothing to delete).
  \item If $X <$ current node's value: recurse left, reattach the
        result.
  \item If $X >$ current node's value: recurse right, reattach the
        result.
  \item Else (found the node to delete):
  \begin{enumerate}
    \item If it has no left child: return its right child (handles both
          the leaf case and the one-right-child case).
    \item If it has no right child: return its left child.
    \item Otherwise (two children): find $\textit{successor} \gets$ the
          minimum value in the right subtree. Set the current node's
          value to $\textit{successor}$. Recursively delete
          $\textit{successor}$ from the right subtree, and reattach the
          result as the new right child.
  \end{enumerate}
  \item Return the (possibly updated) current node.
\end{enumerate}

\subsection*{Dry run}

\dryrun{Delete $X=5$ from the BST with root $5$, left child $3$, right
child $8$ (with left child $7$ and right child $9$).}

At the root ($5=X$): two children present. Find successor: minimum of
right subtree (rooted at $8$) is reached by walking left: $8 \to 7$ (no
left child) -- successor is $7$. Set root's value to $7$. Recursively
delete $7$ from the right subtree (rooted at $8$): at $8$, $7<8$, go
left to $7$; at $7$ ($7=X$ for this recursive call), no children --
return null, attach as $8$'s new left child.

Resulting tree: root $7$, left child $3$, right child $8$ (now with no
left child, right child $9$).

\subsection*{C++ implementation}

\begin{lstlisting}
#include <bits/stdc++.h>
using namespace std;

struct TreeNode {
    int val;
    TreeNode* left;
    TreeNode* right;
    TreeNode(int v) : val(v), left(nullptr), right(nullptr) {}
};

int findMin(TreeNode* root) {
    while (root->left) root = root->left;
    return root->val;
}

TreeNode* deleteNode(TreeNode* root, int X) {
    if (root == nullptr) return nullptr;

    if (X < root->val) {
        root->left = deleteNode(root->left, X);
    } else if (X > root->val) {
        root->right = deleteNode(root->right, X);
    } else {
        if (root->left == nullptr) return root->right;
        if (root->right == nullptr) return root->left;

        int successor = findMin(root->right);
        root->val = successor;
        root->right = deleteNode(root->right, successor);
    }
    return root;
}

int main() {
    TreeNode* root = new TreeNode(5);
    root->left = new TreeNode(3);
    root->right = new TreeNode(8);
    root->right->left = new TreeNode(7);
    root->right->right = new TreeNode(9);

    root = deleteNode(root, 5);
    cout << "New root: " << root->val << endl;
    return 0;
}
\end{lstlisting}

\begin{complexitybox}
\textbf{Time Complexity.} $O(h)$: at most one $O(h)$ walk to find the
node, plus at most one more $O(h)$ walk to find and delete the successor
in the two-children case.
\textbf{Space Complexity.} $O(h)$ for the recursion stack.
\end{complexitybox}

\begin{takeawaybox}
The hard case of BST deletion is always solved the same way: borrow a
boundary value (successor or predecessor) from the subtree that must stay
connected, copy it up, then delete the boundary value's original location
-- which is guaranteed to be an easy case by construction. This
``borrow an edge value, then recursively clean up'' pattern is worth
remembering as its own primitive, independent of the BST context.
\end{takeawaybox}

\section{Floor and Ceil in a BST}
\label{sec:bst-floor-ceil}

\problemstatement{Given the root of a BST and a value $X$, find the
\textbf{floor} of $X$ (the largest value in the tree that is $\le X$) and
the \textbf{ceil} of $X$ (the smallest value in the tree that is $\ge
X$).}

\begin{patternbox}
This is the tree-shaped version of the Binary Search chapter's
Floor and Ceil in a Sorted Array problem: the same ``record a tentative answer, keep narrowing
toward the boundary'' template applies, except narrowing now means moving
to a child rather than shrinking an array window.
\end{patternbox}

\subsection*{Understanding the problem carefully}

Walk from the root, comparing $X$ to each node's value. If the current
node's value is exactly $X$, it is simultaneously its own floor and
ceil -- done immediately. If the current node's value is less than $X$,
it is a \emph{candidate} floor (everything in its left subtree is even
smaller, so it cannot help find a \emph{larger} floor candidate, but the
right subtree might contain an even closer one) -- record it and move
right. Symmetrically, if the current node's value exceeds $X$, it is a
candidate ceil -- record it and move left.

\begin{ideabox}[Record a Candidate, Narrow Toward a Closer One]
Maintain $\textit{floor} \gets -\infty$ (or undefined) and
$\textit{ceil} \gets +\infty$ (or undefined). At each node: if its value
equals $X$, both floor and ceil are this value -- stop. If its value $<
X$: update $\textit{floor}$ to this value (a closer candidate than any
previous one, since we are moving rightward, deeper toward $X$), and move
right. If its value $> X$: update $\textit{ceil}$, and move left.
\end{ideabox}

\subsection*{Why moving right after finding a smaller value always finds a closer floor (if one exists)}

Once we know the current node's value is $<X$ and is therefore a valid
floor candidate, any \emph{better} (larger, but still $\le X$) candidate
can only exist in the current node's right subtree -- the BST invariant
guarantees every value in the right subtree exceeds the current node's
value, narrowing the search toward $X$ from below, while the left subtree
only contains values even smaller, which could never improve on the
current candidate. This mirrors exactly the boundary-narrowing logic of
the Binary Search chapter's array-based floor/ceil, with ``move
right in the tree'' playing the role that ``move the binary search window
toward the boundary'' played there.

\subsection*{Algorithm}

\begin{enumerate}
  \item Initialize $\textit{floor}, \textit{ceil} \gets$ undefined.
  \item While the current node is not null:
  \begin{enumerate}
    \item If current value $=X$: set both $\textit{floor}$ and
          $\textit{ceil}$ to $X$; stop.
    \item If current value $<X$: $\textit{floor} \gets$ current value;
          move to the right child.
    \item Else: $\textit{ceil} \gets$ current value; move to the left
          child.
  \end{enumerate}
  \item Return $(\textit{floor}, \textit{ceil})$.
\end{enumerate}

\subsection*{Dry run}

\dryrun{Take the BST with root $8$, left subtree rooted at $3$ (children
$1$ and $6$, where $6$ has left child $4$), right subtree rooted at
$10$ (left child $9$, right child $14$). Query $X=7$.}

\begin{itemize}
  \item At $8$: $8>7$, $\textit{ceil}=8$. Move left to $3$.
  \item At $3$: $3<7$, $\textit{floor}=3$. Move right to $6$.
  \item At $6$: $6<7$, $\textit{floor}=6$. Move right -- but $6$ has no
        right child (null). Stop.
\end{itemize}

Result: $\textit{floor}=6$, $\textit{ceil}=8$.

\subsection*{C++ implementation}

\begin{lstlisting}
#include <bits/stdc++.h>
using namespace std;

struct TreeNode {
    int val;
    TreeNode* left;
    TreeNode* right;
    TreeNode(int v) : val(v), left(nullptr), right(nullptr) {}
};

pair<int,int> floorCeilBST(TreeNode* root, int X) {
    int floorVal = -1, ceilVal = -1;
    TreeNode* node = root;

    while (node) {
        if (node->val == X) {
            floorVal = ceilVal = X;
            break;
        } else if (node->val < X) {
            floorVal = node->val;
            node = node->right;
        } else {
            ceilVal = node->val;
            node = node->left;
        }
    }
    return {floorVal, ceilVal};
}

int main() {
    TreeNode* root = new TreeNode(8);
    root->left = new TreeNode(3);
    root->left->left = new TreeNode(1);
    root->left->right = new TreeNode(6);
    root->left->right->left = new TreeNode(4);
    root->right = new TreeNode(10);
    root->right->left = new TreeNode(9);
    root->right->right = new TreeNode(14);

    auto [floorVal, ceilVal] = floorCeilBST(root, 7);
    cout << "Floor: " << floorVal << ", Ceil: " << ceilVal << endl;
    return 0;
}
\end{lstlisting}

\begin{complexitybox}
\textbf{Time Complexity.} $O(h)$, a single root-to-(leaf or match) walk.
\textbf{Space Complexity.} $O(1)$ for the iterative version shown (no
recursion needed, since neither subtree is ever revisited).
\end{complexitybox}

\begin{takeawaybox}
Floor and ceil in a BST need only a single $O(h)$ walk, never a full
traversal -- a useful contrast with Kth Smallest
(Section~\ref{sec:bst-kth-smallest}), where, without extra bookkeeping, an
inorder traversal touching many nodes is genuinely required. Whenever a
BST query can be answered by a sequence of single left/right decisions
without needing to inspect both children's full subtrees, prefer the
guided walk over a full traversal.
\end{takeawaybox}

\section{Kth Smallest Element in a BST}
\label{sec:bst-kth-smallest}

\problemstatement{Given the root of a BST and an integer $k$, find the
$k$-th smallest value stored in the tree.}

\begin{patternbox}
\emph{``Kth smallest''} on a BST is answered directly by
Section~\ref{sec:bst-intro}'s headline fact: an \textbf{inorder
traversal visits values in strictly increasing order}. So the $k$-th
value produced by an inorder traversal \emph{is} the $k$-th smallest
value in the tree -- no sorting, and no need to visit the whole tree if
we stop as soon as the $k$-th value is produced.
\end{patternbox}

\subsection*{Understanding the problem carefully}

A recursive inorder traversal could be used, but stopping it early (the
moment the $k$-th value is produced) is awkward with plain recursion,
since recursive calls don't have a natural way to ``stop everything''
once we are several levels deep. An \textbf{iterative} inorder traversal,
using an explicit stack to simulate the recursion, gives full control over
exactly when to stop -- we simply count nodes as we pop them and halt the
moment the count reaches $k$.

\begin{ideabox}[Iterative Inorder Traversal, Counting as We Go]
Maintain an explicit stack. Repeatedly: push the current node and move
left, until reaching null (this pushes the entire leftward chain, since
inorder visits left subtrees first). Then pop a node -- this is the next
value in sorted order; increment a counter, and if the counter reaches
$k$, return this node's value immediately. Otherwise, move to the popped
node's right child and repeat the whole process (push its leftward
chain, then continue popping).
\end{ideabox}

\subsection*{Why the iterative version visits nodes in exactly inorder sequence}

Pushing the entire leftward chain before popping anything ensures that,
when we do pop, we are popping the smallest value not yet visited: any
node still left to visit that is smaller than the current pop candidate
would have to be in some left subtree along the path already pushed, and
all of those are already on the stack, ready to be popped first (stacks
pop most-recently-pushed first, which corresponds to the deepest,
leftmost, smallest unvisited node). After popping a node, moving to its
right child and repeating correctly resumes the traversal exactly where
recursive inorder would have continued, since everything in that node's
own left subtree has already been fully visited.

\subsection*{Algorithm}

\begin{enumerate}
  \item Initialize an empty stack, $\textit{count} \gets 0$,
        $\textit{node} \gets \textit{root}$.
  \item While $\textit{node}$ is not null or the stack is not empty:
  \begin{enumerate}
    \item While $\textit{node}$ is not null: push it; $\textit{node}
          \gets \textit{node}.\textit{left}$.
    \item Pop $\textit{node}$; $\textit{count} \gets \textit{count}+1$.
    \item If $\textit{count} = k$: return $\textit{node}.\textit{val}$.
    \item $\textit{node} \gets \textit{node}.\textit{right}$.
  \end{enumerate}
\end{enumerate}

\subsection*{Dry run}

\dryrun{Take the BST with root $5$, left child $3$ (children $1,4$),
right child $8$, $k=3$.}

\begin{itemize}
  \item Push leftward chain from $5$: push $5$, push $3$, push $1$ (no
        left child). Stack (top to bottom) $=[1,3,5]$.
  \item Pop $1$: count$=1 \ne 3$. Move to $1$'s right (null).
  \item Push leftward chain from null: nothing to push.
  \item Pop $3$: count$=2 \ne 3$. Move to $3$'s right, which is $4$.
  \item Push leftward chain from $4$: push $4$ (no left child).
  \item Pop $4$: count$=3=k$. Return $4$.
\end{itemize}

The $3$rd smallest value is $4$ (inorder sequence: $1,3,4,5,8$).

\subsection*{C++ implementation}

\begin{lstlisting}
#include <bits/stdc++.h>
using namespace std;

struct TreeNode {
    int val;
    TreeNode* left;
    TreeNode* right;
    TreeNode(int v) : val(v), left(nullptr), right(nullptr) {}
};

int kthSmallest(TreeNode* root, int k) {
    stack<TreeNode*> st;
    TreeNode* node = root;
    int count = 0;

    while (node || !st.empty()) {
        while (node) {
            st.push(node);
            node = node->left;
        }
        node = st.top();
        st.pop();
        count++;
        if (count == k) return node->val;
        node = node->right;
    }
    return -1; // k larger than tree size
}

int main() {
    TreeNode* root = new TreeNode(5);
    root->left = new TreeNode(3);
    root->right = new TreeNode(8);
    root->left->left = new TreeNode(1);
    root->left->right = new TreeNode(4);

    cout << "3rd smallest: " << kthSmallest(root, 3) << endl;
    return 0;
}
\end{lstlisting}

\begin{complexitybox}
\textbf{Time Complexity.} $O(h+k)$: pushing the initial leftward chain
costs $O(h)$, and each subsequent step processes one more node until the
$k$-th is found, giving $O(h+k)$ overall -- in the worst case $O(n)$ if
$k$ is close to the total node count.
\textbf{Space Complexity.} $O(h)$ for the stack.
\end{complexitybox}

\begin{takeawaybox}
Whenever a BST problem needs the $i$-th value in sorted order, or values
processed one at a time in sorted order with the ability to stop early,
an \emph{iterative} inorder traversal with an explicit stack is the right
tool -- it gives the same sorted-order guarantee as recursive inorder,
with full control over when to stop. This exact iterative inorder
machinery, generalized to resume across multiple separate calls, is what
Section~\ref{sec:bst-iterator}'s BST Iterator formalizes next.
\end{takeawaybox}

\section{Validate a Binary Search Tree}
\label{sec:bst-validate}

\problemstatement{Given the root of a binary tree, determine whether it
is a valid BST: for every node, all values in its left subtree must be
strictly less than the node's value, and all values in its right subtree
must be strictly greater.}

\begin{patternbox}
The keyword trap here: it is tempting to check only that each node is
greater than its immediate left child and less than its immediate right
child, but the BST invariant is actually about \textbf{every} node in
each subtree, not just the direct children. \emph{``Validate a property
that must hold against \textbf{all ancestors}, not just the immediate
parent''} signals that the recursion needs to carry a \textbf{valid
range} downward, narrowing at every step, rather than just comparing
parent to child.
\end{patternbox}

\subsection*{Understanding the problem carefully}

Consider the tree with root $5$, left child $3$, and $3$'s right child
$6$. Checking only parent-child pairs would see: $3<5$ (fine, left child
smaller than parent) and $6>3$ (fine, right child larger than its own
parent) -- and incorrectly conclude this is a valid BST. But $6$ sits in
$5$'s \emph{left} subtree, and the BST invariant requires \emph{every}
value there to be less than $5$, which $6$ violates. The bug is checking
only local parent-child relationships instead of the full inherited
range of validity.

\begin{ideabox}[Pass Down a Valid (Lower, Upper) Range]
Recurse with two bounds, $\textit{lower}$ and $\textit{upper}$
(initially $-\infty$ and $+\infty$ at the root). At each node: if its
value is not strictly between $\textit{lower}$ and $\textit{upper}$,
the tree is invalid. Otherwise, recurse left with the same
$\textit{lower}$ but $\textit{upper}$ tightened to the current node's
value (everything in the left subtree must be less than this node), and
recurse right with $\textit{lower}$ tightened to the current node's
value and the same $\textit{upper}$.
\end{ideabox}

\subsection*{Why narrowing the range at every level correctly enforces the global invariant}

By construction, the range $(\textit{lower}, \textit{upper})$ passed into
any node already reflects the combined constraints of \emph{every}
ancestor on the path from the root, not just the immediate parent: each
recursive call only ever tightens the range, never loosens it, so a deep
descendant's range is the \emph{intersection} of every constraint imposed
by every ancestor above it. Checking the current node's value against
this fully-accumulated range is therefore equivalent to checking it
against every single ancestor directly, just done incrementally and in
$O(1)$ work per node instead of re-walking ancestors at every step.

\subsection*{Algorithm}

\begin{enumerate}
  \item \textsc{Validate}(\textit{node}, \textit{lower}, \textit{upper}):
  \begin{enumerate}
    \item If $\textit{node} = $ null: return true (an empty tree is
          trivially valid).
    \item If $\textit{node}.\textit{val} \le \textit{lower}$ or
          $\textit{node}.\textit{val} \ge \textit{upper}$: return
          false.
    \item Return
          \textsc{Validate}($\textit{node}.\textit{left}, \textit{lower},
          \textit{node}.\textit{val}$) \textbf{and}
          \textsc{Validate}($\textit{node}.\textit{right},
          \textit{node}.\textit{val}, \textit{upper}$).
  \end{enumerate}
  \item Call with $\textit{lower}=-\infty$, $\textit{upper}=+\infty$ at
        the root.
\end{enumerate}

\subsection*{Dry run}

\dryrun{Take the tree with root $5$, left child $3$, $3$'s right child
$6$.}

\begin{itemize}
  \item At $5$, range $(-\infty,+\infty)$: $5$ is within range.
        Recurse left into $3$ with range $(-\infty,5)$; recurse right
        into null with range $(5,+\infty)$ (trivially valid).
  \item At $3$, range $(-\infty,5)$: $3$ is within range. Recurse left
        into null (valid); recurse right into $6$ with range $(3,5)$.
  \item At $6$, range $(3,5)$: is $6$ strictly between $3$ and $5$? No
        ($6 \ge 5$). Return false.
\end{itemize}

The tree is correctly identified as \textbf{invalid}.

\subsection*{C++ implementation}

\begin{lstlisting}
#include <bits/stdc++.h>
using namespace std;

struct TreeNode {
    int val;
    TreeNode* left;
    TreeNode* right;
    TreeNode(int v) : val(v), left(nullptr), right(nullptr) {}
};

bool validate(TreeNode* node, long long lower, long long upper) {
    if (node == nullptr) return true;
    if (node->val <= lower || node->val >= upper) return false;

    return validate(node->left, lower, node->val) &&
           validate(node->right, node->val, upper);
}

bool isValidBST(TreeNode* root) {
    return validate(root, LLONG_MIN, LLONG_MAX);
}

int main() {
    TreeNode* root = new TreeNode(5);
    root->left = new TreeNode(3);
    root->left->right = new TreeNode(6);

    cout << (isValidBST(root) ? "true" : "false") << endl;
    return 0;
}
\end{lstlisting}

\begin{complexitybox}
\textbf{Time Complexity.} Every node is visited exactly once, with
$O(1)$ work per node, giving
\[
  O(n).
\]
\textbf{Space Complexity.} $O(h)$ for the recursion stack.
\end{complexitybox}

\begin{takeawaybox}
Any property phrased as ``true for every node relative to \emph{all} of
its ancestors/descendants,'' not just its immediate neighbors, needs
state (here, a valid range) threaded through the recursion, accumulating
constraints as it descends. An equally valid alternative solution
performs an inorder traversal and checks that the resulting sequence is
strictly increasing, directly invoking
Section~\ref{sec:bst-intro}'s headline fact -- worth knowing as a second
way to arrive at the same answer.
\end{takeawaybox}

\section{Predecessor and Successor in a BST}
\label{sec:bst-pred-succ}

\problemstatement{Given the root of a BST and a key $X$ (which may or may
not itself be present in the tree), find the \textbf{inorder predecessor}
of $X$ (the largest value in the tree strictly less than $X$) and the
\textbf{inorder successor} of $X$ (the smallest value strictly greater
than $X$).}

\begin{patternbox}
Predecessor and successor are floor and ceil
(Section~\ref{sec:bst-floor-ceil}) with \textbf{strict} inequalities
instead of non-strict ones -- the keyword difference between ``floor''
(largest $\le X$) and ``predecessor'' (largest $<X$) is exactly the
difference between $\ge$ and $>$ that distinguished lower bound from
upper bound back in the Binary Search chapter.
\end{patternbox}

\subsection*{Understanding the problem carefully}

The walk is identical in shape to floor/ceil: at each node, compare to
$X$ and decide which candidate to update and which direction to move.
The only change is that when the current node's value equals $X$
\emph{exactly}, we must \textbf{not} stop immediately (unlike floor/ceil,
where an exact match ends the search for both) -- a predecessor or
successor is required to be \emph{strictly} different from $X$, so if the
current node's value equals $X$, we must continue descending: moving left
to find a predecessor strictly smaller, or right to find a successor
strictly larger.

\begin{ideabox}[Same Walk as Floor/Ceil, but Treat an Exact Match as ``Keep Going'']
For the predecessor: at each node, if its value $<X$ (this includes the
case of looking for a value strictly less, so equality to $X$ does
\emph{not} qualify): record it as the current best predecessor candidate,
and move right (search for an even closer one). If its value $\ge X$: do
not record anything; move left. Symmetrically for the successor, with
roles of $<$/$\ge$ swapped to $>$/$\le$.
\end{ideabox}

\subsection*{Why treating equality as ``move past it'' (not ``stop'') is the only change needed}

In floor/ceil, an exact match is itself simultaneously the floor and the
ceil, so the search correctly terminates there. In predecessor/successor,
an exact match is explicitly excluded from being its own answer, so the
search must instead behave exactly as if that node were just slightly
too small (for predecessor purposes, treat it like any other ``$\ge X$''
node and move left without recording it) or slightly too large (for
successor purposes, move right without recording it) -- everything else
about the boundary-narrowing logic carries over unchanged from
Section~\ref{sec:bst-floor-ceil}.

\subsection*{Algorithm}

\begin{enumerate}
  \item \textsc{Predecessor}: $\textit{pred} \gets$ undefined,
        $\textit{node} \gets \textit{root}$. While $\textit{node}$ is
        not null: if $\textit{node}.\textit{val} < X$:
        $\textit{pred} \gets \textit{node}.\textit{val}$; move right.
        Else: move left.
  \item \textsc{Successor}: symmetric, with $>$ in place of $<$ and
        left/right swapped.
\end{enumerate}

\subsection*{Dry run}

\dryrun{Take the BST with root $8$, left child $3$ (children $1,6$, where
$6$ has left child $4$), right child $10$ (children $9,14$). Query
$X=6$.}

\textbf{Predecessor:} At $8$: $8<6$? No, move left to $3$. At $3$:
$3<6$? Yes, $\textit{pred}=3$, move right to $6$. At $6$: $6<6$? No
(equal, not strictly less), move left to $4$. At $4$: $4<6$? Yes,
$\textit{pred}=4$, move right (null). Stop. Predecessor $=4$.

\textbf{Successor:} At $8$: $8>6$? Yes, $\textit{succ}=8$, move left to
$3$. At $3$: $3>6$? No, move right to $6$. At $6$: $6>6$? No (equal),
move right (null). Stop. Successor $=8$.

\subsection*{C++ implementation}

\begin{lstlisting}
#include <bits/stdc++.h>
using namespace std;

struct TreeNode {
    int val;
    TreeNode* left;
    TreeNode* right;
    TreeNode(int v) : val(v), left(nullptr), right(nullptr) {}
};

int findPredecessor(TreeNode* root, int X) {
    int pred = -1;
    while (root) {
        if (root->val < X) {
            pred = root->val;
            root = root->right;
        } else {
            root = root->left;
        }
    }
    return pred;
}

int findSuccessor(TreeNode* root, int X) {
    int succ = -1;
    while (root) {
        if (root->val > X) {
            succ = root->val;
            root = root->left;
        } else {
            root = root->right;
        }
    }
    return succ;
}

int main() {
    TreeNode* root = new TreeNode(8);
    root->left = new TreeNode(3);
    root->left->left = new TreeNode(1);
    root->left->right = new TreeNode(6);
    root->left->right->left = new TreeNode(4);
    root->right = new TreeNode(10);
    root->right->left = new TreeNode(9);
    root->right->right = new TreeNode(14);

    cout << "Predecessor: " << findPredecessor(root, 6) << endl;
    cout << "Successor: " << findSuccessor(root, 6) << endl;
    return 0;
}
\end{lstlisting}

\begin{complexitybox}
\textbf{Time Complexity.} $O(h)$ for each of predecessor and successor.
\textbf{Space Complexity.} $O(1)$ for the iterative version shown.
\end{complexitybox}

\begin{takeawaybox}
Predecessor/successor and floor/ceil are the same boundary-search
template; the only behavioral difference is what happens on an exact
match (stop immediately for floor/ceil, keep narrowing for
predecessor/successor). When given a problem statement, check carefully
whether equality should count as a valid answer or must be strictly
excluded -- that single detail is the entire distinction between these
two near-identical problems.
\end{takeawaybox}

\section{BST Iterator}
\label{sec:bst-iterator}

\problemstatement{Design an iterator over a BST that supports
$\textit{next}()$ (returns the next smallest value in the tree, in
sorted order) and $\textit{hasNext}()$ (whether any values remain),
initialized with the tree's root. Across the iterator's whole lifetime,
$\textit{next}$ should average $O(1)$ time and the iterator should use
$O(h)$ space, not $O(n)$.}

\begin{patternbox}
\emph{``Iterator returning sorted values one at a time''} is exactly
Section~\ref{sec:bst-kth-smallest}'s iterative inorder traversal,
\textbf{paused and resumed} across separate method calls instead of run
to completion inside a single function. The keyword shift from
``compute the $k$-th value'' to ``design an iterator'' is a signal that
the same stack-based traversal state needs to persist as object state
between calls, rather than living entirely within one function's local
variables.
\end{patternbox}

\subsection*{Understanding the problem carefully}

If we simply ran a full inorder traversal up front and stored all $n$
values in a list, $\textit{next}()$ would be trivially $O(1)$, but
construction would cost $O(n)$ time \emph{and} space -- failing the
$O(h)$ space requirement. The fix is to keep \emph{only} the explicit
stack from Section~\ref{sec:bst-kth-smallest}'s iterative traversal alive
as the iterator's state, advancing it by exactly one step per
$\textit{next}()$ call instead of running it to completion all at once.

\begin{ideabox}[Persist the Traversal Stack as Iterator State]
On construction: push the entire leftward chain from the root onto the
stack (exactly as in Section~\ref{sec:bst-kth-smallest}'s setup).
$\textit{hasNext}()$: simply check whether the stack is non-empty.
$\textit{next}()$: pop the top of the stack (this is the next value in
sorted order); then push the leftward chain starting from the popped
node's right child (preparing the stack so that the \emph{following}
call to \textit{next} will correctly find the next value after this
one); return the popped value.
\end{ideabox}

\subsection*{Why pushing only the new leftward chain after each pop keeps amortized cost low}

Each node in the tree is pushed onto the stack exactly once during the
iterator's entire lifetime (either during construction, for nodes on the
initial leftward chain, or during some later $\textit{next}()$ call,
when it becomes the leftward-chain push triggered by its parent being
popped), and each node is popped exactly once. So the \emph{total} work
across all calls to $\textit{next}()$ over the iterator's lifetime is
$O(n)$ (each node contributes $O(1)$ amortized push/pop work) -- which,
spread over the $n$ calls needed to exhaust the tree, gives $O(1)$
\emph{average} time per call, even though any single call might cost up
to $O(h)$ if it happens to push a long leftward chain.

\subsection*{Algorithm}

\begin{enumerate}
  \item \textsc{Constructor}(\textit{root}): push the leftward chain
        from \textit{root} onto the stack.
  \item \textsc{HasNext}(): return whether the stack is non-empty.
  \item \textsc{Next}(): pop $\textit{node}$ from the stack; push the
        leftward chain starting from $\textit{node}.\textit{right}$;
        return $\textit{node}.\textit{val}$.
\end{enumerate}

\subsection*{Dry run}

\dryrun{Take the BST with root $7$, left child $3$, right child $15$
(left child $9$, right child $20$). Call \textit{next} three times.}

\begin{itemize}
  \item Construct: push leftward chain from $7$: push $7$, push $3$ (no
        left child). Stack (top to bottom) $=[3,7]$.
  \item \textit{next}(): pop $3$. Push leftward chain from $3$'s right
        (null) -- nothing to push. Stack $=[7]$. Return $3$.
  \item \textit{next}(): pop $7$. Push leftward chain from $7$'s right,
        $15$: push $15$, then push $15$'s left, $9$ (no left child).
        Stack $=[9,15]$. Return $7$.
  \item \textit{next}(): pop $9$. Push leftward chain from $9$'s right
        (null) -- nothing. Stack $=[15]$. Return $9$.
\end{itemize}

Output sequence so far: $3,7,9$ -- correctly sorted order.

\subsection*{C++ implementation}

\begin{lstlisting}
#include <bits/stdc++.h>
using namespace std;

struct TreeNode {
    int val;
    TreeNode* left;
    TreeNode* right;
    TreeNode(int v) : val(v), left(nullptr), right(nullptr) {}
};

class BSTIterator {
    stack<TreeNode*> st;

    void pushLeftChain(TreeNode* node) {
        while (node) {
            st.push(node);
            node = node->left;
        }
    }

public:
    BSTIterator(TreeNode* root) {
        pushLeftChain(root);
    }

    bool hasNext() {
        return !st.empty();
    }

    int next() {
        TreeNode* node = st.top();
        st.pop();
        pushLeftChain(node->right);
        return node->val;
    }
};

int main() {
    TreeNode* root = new TreeNode(7);
    root->left = new TreeNode(3);
    root->right = new TreeNode(15);
    root->right->left = new TreeNode(9);
    root->right->right = new TreeNode(20);

    BSTIterator it(root);
    while (it.hasNext()) cout << it.next() << " ";
    cout << endl;
    return 0;
}
\end{lstlisting}

\begin{complexitybox}
\textbf{Time Complexity.} $O(1)$ average per $\textit{next}()$ call
(amortized over the iterator's lifetime), $O(h)$ worst case for any
single call; $\textit{hasNext}()$ is always $O(1)$.
\textbf{Space Complexity.} $O(h)$ for the stack, never $O(n)$.
\end{complexitybox}

\begin{takeawaybox}
Converting a one-shot traversal into a resumable iterator is purely a
matter of moving the traversal's local state (here, the explicit stack)
into object-level state that persists between calls. The traversal logic
itself -- push the leftward chain, pop, push the popped node's right
subtree's leftward chain -- is completely unchanged from
Section~\ref{sec:bst-kth-smallest}.
\end{takeawaybox}

\section{Two Sum IV: Input is a BST}
\label{sec:bst-two-sum}

\problemstatement{Given the root of a BST and a target integer $k$,
determine whether there exist two \textbf{distinct} nodes whose values
sum to exactly $k$.}

\begin{patternbox}
\emph{``Two sum''} on a BST is the classic two-pointer two-sum technique
(normally run on a sorted array) combined with
Section~\ref{sec:bst-iterator}'s BST Iterator: since an inorder traversal
already produces values in sorted order on demand, we can drive
\textbf{two} iterators -- one walking forward (ascending) from the
smallest value, one walking backward (descending) from the largest --
exactly as the classic two-pointer two-sum algorithm drives two pointers
inward from the two ends of a sorted array, without ever materializing
that sorted array.
\end{patternbox}

\subsection*{Understanding the problem carefully}

If the tree's values were laid out as a sorted array, the standard
two-sum technique would maintain a pointer at each end: compute the sum
of the two pointed-at values; if it equals $k$, done; if it is too
small, advance the left pointer rightward (to a larger value); if too
large, move the right pointer leftward (to a smaller value). A BST
supports exactly this access pattern \emph{without} an explicit array,
using two independent, resumable inorder-traversal iterators: one
producing ascending values on demand (Section~\ref{sec:bst-iterator}'s
iterator exactly as built), and a mirror-image iterator producing
descending values on demand (the same idea, with left/right swapped).

\begin{ideabox}[Two BST Iterators, Driven Like Two-Sum's Two Pointers]
Initialize a forward iterator $\textit{lo}$ (smallest-first) and a
backward iterator $\textit{hi}$ (largest-first), both via the stack-based
technique from Section~\ref{sec:bst-iterator}. Maintain their current
values $a = \textit{lo}.\textit{next}()$ and $b =
\textit{hi}.\textit{next}()$. While $a < b$: if $a+b=k$, return true; if
$a+b<k$, advance $\textit{lo}$ (get a larger $a$); if $a+b>k$, advance
$\textit{hi}$ (get a smaller $b$). If the pointers meet or cross ($a \ge
b$) without finding a match, return false.
\end{ideabox}

\subsection*{Why driving two independent stack-based iterators is correct and efficient}

This is exactly the two-pointer two-sum algorithm, with the ``array'' it
normally operates on replaced by two on-demand sorted streams generated
directly from the tree. Each iterator independently maintains its own
$O(h)$-space stack, advancing forward or backward through the BST's
sorted order exactly as Section~\ref{sec:bst-iterator} established, and
because both iterators only ever move \emph{toward} each other (never
revisiting a value), the total work across the entire search is bounded
by the number of \emph{distinct} values each iterator visits before they
cross -- giving the same linear-total-work guarantee as the classic
array-based two-pointer technique, without ever needing $O(n)$ space to
store a flattened sorted array.

\subsection*{Algorithm}

\begin{enumerate}
  \item Initialize a forward BST iterator $\textit{lo}$ and a backward
        BST iterator $\textit{hi}$.
  \item $a \gets \textit{lo}.\textit{next}()$, $b \gets
        \textit{hi}.\textit{next}()$.
  \item While $a < b$:
  \begin{enumerate}
    \item If $a+b=k$: return true.
    \item If $a+b<k$: $a \gets \textit{lo}.\textit{next}()$.
    \item Else: $b \gets \textit{hi}.\textit{next}()$.
  \end{enumerate}
  \item Return false.
\end{enumerate}

\subsection*{Dry run}

\dryrun{Take the BST with root $5$, left child $3$, right child $8$,
target $k=10$.}

Forward iterator yields $3,5,8$ in order; backward iterator yields
$8,5,3$ in order.

\begin{itemize}
  \item $a=\textit{lo.next}()=3$, $b=\textit{hi.next}()=8$. $a<b$.
        $a+b=11>10$. Advance \textit{hi}: $b=5$.
  \item $a=3<b=5$. $a+b=8<10$. Advance \textit{lo}: $a=5$.
  \item $a=5$, $b=5$: $a<b$ is now false ($5 \not< 5$). Loop ends.
\end{itemize}

Return false. Checking by hand: the only pair sums available are
$3{+}5{=}8$, $3{+}8{=}11$, $5{+}8{=}13$ -- none equal $10$. Correct.

\subsection*{C++ implementation}

\begin{lstlisting}
#include <bits/stdc++.h>
using namespace std;

struct TreeNode {
    int val;
    TreeNode* left;
    TreeNode* right;
    TreeNode(int v) : val(v), left(nullptr), right(nullptr) {}
};

class BSTIterator {
    stack<TreeNode*> st;
    bool reverseOrder; // true = descending (successor uses right first)

    void pushChain(TreeNode* node) {
        while (node) {
            st.push(node);
            node = reverseOrder ? node->right : node->left;
        }
    }

public:
    BSTIterator(TreeNode* root, bool reverse) : reverseOrder(reverse) {
        pushChain(root);
    }

    int next() {
        TreeNode* node = st.top();
        st.pop();
        pushChain(reverseOrder ? node->left : node->right);
        return node->val;
    }
};

bool findTarget(TreeNode* root, int k) {
    BSTIterator lo(root, false); // ascending
    BSTIterator hi(root, true);  // descending

    int a = lo.next(), b = hi.next();
    while (a < b) {
        if (a + b == k) return true;
        else if (a + b < k) a = lo.next();
        else b = hi.next();
    }
    return false;
}

int main() {
    TreeNode* root = new TreeNode(5);
    root->left = new TreeNode(3);
    root->right = new TreeNode(8);

    cout << (findTarget(root, 10) ? "true" : "false") << endl;
    cout << (findTarget(root, 11) ? "true" : "false") << endl;
    return 0;
}
\end{lstlisting}

\begin{complexitybox}
\textbf{Time Complexity.} $O(n)$ in the worst case -- each value is
produced by one of the two iterators at most once before they cross,
and each iterator's amortized cost per call is $O(1)$, as established in
Section~\ref{sec:bst-iterator}.
\textbf{Space Complexity.} $O(h)$ for each iterator's stack, so $O(h)$
total.
\end{complexitybox}

\begin{takeawaybox}
A BST is, for the purposes of any algorithm that only needs sorted-order
access, a drop-in replacement for a sorted array -- as long as you have
an $O(h)$-space, resumable iterator (forward and/or backward) instead of
an $O(n)$-space materialized array. Any two-pointer array technique
(two-sum, partitioning, merging) can be ported to a BST this way, exactly
as done here.
\end{takeawaybox}

\section{Lowest Common Ancestor in a BST}
\label{sec:bst-lca}

\problemstatement{Given the root of a BST and two nodes $p$ and $q$
(both guaranteed to exist in the tree), find their \textbf{lowest common
ancestor} (LCA): the deepest node that has both $p$ and $q$ as
descendants (a node is considered a descendant of itself).}

\begin{patternbox}
LCA in a general binary tree requires examining both subtrees at every
node and combining results bottom-up, costing $O(n)$. The keyword
\emph{``BST''} attached to an LCA problem is a strong signal that the
BST invariant alone determines the split point directly, top-down, in
$O(h)$ -- no need to inspect both subtrees blindly the way a general
tree's LCA algorithm would.
\end{patternbox}

\subsection*{Understanding the problem carefully}

In a BST, the LCA of $p$ and $q$ is exactly the node where the search
paths to $p$ and $q$ \textbf{diverge}: as long as both $p.\textit{val}$
and $q.\textit{val}$ are on the same side (both less than the current
node, or both greater), the current node cannot be the LCA -- both
targets are reachable only by continuing in that single, same direction,
so the LCA must lie further down on that side. The moment $p$ and $q$
fall on \emph{different} sides (or one of them \emph{is} the current
node), the current node is exactly the point where their paths separate
-- the LCA.

\begin{ideabox}[Follow Both Targets Together Until They Diverge]
Starting at the root: if both $p.\textit{val}$ and $q.\textit{val}$ are
less than the current node's value, move left. If both are greater, move
right. Otherwise (one is $\le$ and the other is $\ge$ the current node's
value, including the case where the current node \emph{is} $p$ or $q$),
the current node is the LCA.
\end{ideabox}

\subsection*{Why the divergence point is provably the LCA}

If both targets are smaller than the current node, the current node
cannot be an ancestor that is also the \emph{lowest} one -- both $p$ and
$q$ are confirmed to live in the left subtree, so a deeper common
ancestor must exist there, and continuing left can only get closer to it
(never skip past it, since the BST invariant guarantees the entire left
subtree is exactly where both targets are). The symmetric argument holds
for both being larger. The first node where this is no longer true --
where the two targets are no longer both on the same side -- is, by
definition, the deepest point both search paths still share, since any
node below it would have to continue toward only one of $p$ or $q$, not
both.

\subsection*{Algorithm}

\begin{enumerate}
  \item Start at the root.
  \item While true:
  \begin{enumerate}
    \item If $p.\textit{val} < \textit{current}.\textit{val}$ and
          $q.\textit{val} < \textit{current}.\textit{val}$: move to the
          left child.
    \item Else if $p.\textit{val} > \textit{current}.\textit{val}$ and
          $q.\textit{val} > \textit{current}.\textit{val}$: move to the
          right child.
    \item Else: return \textit{current} (this is the LCA).
  \end{enumerate}
\end{enumerate}

\subsection*{Dry run}

\dryrun{Take the BST with root $6$, left child $2$ (children $0,4$),
right child $8$ (children $7,9$). Find LCA of $p=2$, $q=4$.}

\begin{itemize}
  \item At $6$: is $2<6$ and $4<6$? Yes, both smaller. Move left to
        $2$.
  \item At $2$: is $2<2$? No ($p$ equals the current node). So not both
        strictly smaller, and not both strictly larger. Return $2$.
\end{itemize}

LCA is $2$ (correct: $2$ is an ancestor of $4$, and is itself $p$).

\subsection*{C++ implementation}

\begin{lstlisting}
#include <bits/stdc++.h>
using namespace std;

struct TreeNode {
    int val;
    TreeNode* left;
    TreeNode* right;
    TreeNode(int v) : val(v), left(nullptr), right(nullptr) {}
};

TreeNode* lowestCommonAncestor(TreeNode* root, TreeNode* p, TreeNode* q) {
    TreeNode* current = root;
    while (true) {
        if (p->val < current->val && q->val < current->val) {
            current = current->left;
        } else if (p->val > current->val && q->val > current->val) {
            current = current->right;
        } else {
            return current;
        }
    }
}

int main() {
    TreeNode* root = new TreeNode(6);
    root->left = new TreeNode(2);
    root->right = new TreeNode(8);
    root->left->left = new TreeNode(0);
    root->left->right = new TreeNode(4);
    root->right->left = new TreeNode(7);
    root->right->right = new TreeNode(9);

    TreeNode* lca = lowestCommonAncestor(root, root->left, root->left->right);
    cout << "LCA: " << lca->val << endl;
    return 0;
}
\end{lstlisting}

\begin{complexitybox}
\textbf{Time Complexity.} $O(h)$ -- a single root-to-LCA walk, never
touching any subtree that doesn't lie on the path to the answer.
\textbf{Space Complexity.} $O(1)$ for the iterative version shown.
\end{complexitybox}

\begin{takeawaybox}
LCA in a general binary tree needs $O(n)$ because there is no way to know
which subtree contains which target without searching both. The BST
invariant removes this ambiguity entirely: comparing both targets'
values to the current node tells you, with certainty, which single
direction to go (or that you have already arrived), collapsing an
$O(n)$ two-subtree search into an $O(h)$ single-path walk.
\end{takeawaybox}

\section{Construct a BST from Preorder Traversal}
\label{sec:bst-construct-preorder}

\problemstatement{Given an array representing the preorder traversal of a
BST (root, then left subtree, then right subtree, with no duplicate
values), reconstruct the BST and return its root.}

\begin{patternbox}
\emph{``Reconstruct a BST from one traversal''} works only because the
traversal is specifically \textbf{preorder} \emph{and} the tree is
specifically a \textbf{BST} -- the very first element is always the root,
and the BST invariant (combined with knowing each value's relationship to
the root) is enough to tell us, for every subsequent element, whether it
belongs in the left subtree or the right, without needing a second
traversal (like inorder) to disambiguate, the way general binary tree
reconstruction usually requires.
\end{patternbox}

\subsection*{Understanding the problem carefully}

In a preorder traversal of a BST, the first value is the root. Every
value immediately following it that is \textbf{less than} the root
belongs to the root's \emph{left} subtree -- and because of how preorder
works (left subtree fully traversed before any of the right subtree
begins), \emph{all} such smaller values appear contiguously, immediately
after the root, before any larger value appears. The first value that is
\textbf{greater} than the root marks the start of the right subtree, and
every later value will also be greater than the root.

\begin{ideabox}[Build with an Upper Bound, Consuming a Shared Pointer]
Maintain a single index pointer into the preorder array, shared across
all recursive calls (advanced every time a node is consumed). Define
$\textsc{Build}(\textit{upperBound})$: if the pointer is exhausted, or
the next value to consume exceeds $\textit{upperBound}$, return null
(this value belongs to some ancestor's other subtree, not here). Otherwise,
consume the next value as a new node; recursively build its left subtree
with $\textit{upperBound}$ tightened to this new node's value; then
recursively build its right subtree with the \emph{original}
$\textit{upperBound}$ unchanged (everything remaining that is still
$\le$ the original bound, but now also $>$ this node's value, belongs on
the right).
\end{ideabox}

\subsection*{Why a single shared pointer and a bound (not an explicit split search) suffices}

The key insight is that we never need to explicitly search for ``where
the left subtree ends and the right subtree begins'' -- the recursive
calls naturally stop consuming values the moment they encounter one that
violates the current $\textit{upperBound}$, which is exactly the boundary
between subtrees. Because the shared pointer only ever advances (never
resets or searches backward), each value in the preorder array is
considered, and consumed, exactly once across the entire construction,
regardless of how deeply nested its eventual position in the tree is.

\subsection*{Algorithm}

\begin{enumerate}
  \item Initialize a shared index $i \gets 0$.
  \item $\textsc{Build}(\textit{upperBound})$:
  \begin{enumerate}
    \item If $i \ge n$ or $\textit{preorder}[i] > \textit{upperBound}$:
          return null.
    \item Create a new node with value $\textit{preorder}[i]$;
          increment $i$.
    \item Recursively set this node's left child to
          $\textsc{Build}(\textit{node.val})$.
    \item Recursively set this node's right child to
          $\textsc{Build}(\textit{upperBound})$ (the original bound).
    \item Return this node.
  \end{enumerate}
  \item Call $\textsc{Build}(+\infty)$ at the top level.
\end{enumerate}

\subsection*{Dry run}

\dryrun{Take $\textit{preorder} = [8,5,1,7,10,12]$.}

\begin{itemize}
  \item $\textsc{Build}(+\infty)$: consume $8$ ($i \to 1$). Node $8$.
  \item Left: $\textsc{Build}(8)$: consume $5$ ($5\le8$, $i\to2$). Node
        $5$.
  \begin{itemize}
    \item Left: $\textsc{Build}(5)$: consume $1$ ($1\le5$, $i\to3$).
          Node $1$. Left: $\textsc{Build}(1)$: next is $7>1$, return
          null. Right: $\textsc{Build}(5)$: next is $7>5$, return null.
          Node $1$ is a leaf.
    \item Right: $\textsc{Build}(8)$: consume $7$ ($7\le8$, $i\to4$).
          Node $7$. Left: $\textsc{Build}(7)$: next is $10>7$, return
          null. Right: $\textsc{Build}(8)$: next is $10>8$, return null.
          Node $7$ is a leaf.
  \end{itemize}
  \item Right of $8$: $\textsc{Build}(+\infty)$: consume $10$
        ($i\to5$). Node $10$.
  \begin{itemize}
    \item Left: $\textsc{Build}(10)$: next is $12>10$, return null.
    \item Right: $\textsc{Build}(+\infty)$: consume $12$ ($i\to6$).
          Node $12$, a leaf.
  \end{itemize}
\end{itemize}

Resulting tree: root $8$, left child $5$ (children $1,7$), right child
$10$ (right child $12$).

\subsection*{C++ implementation}

\begin{lstlisting}
#include <bits/stdc++.h>
using namespace std;

struct TreeNode {
    int val;
    TreeNode* left;
    TreeNode* right;
    TreeNode(int v) : val(v), left(nullptr), right(nullptr) {}
};

TreeNode* build(vector<int>& preorder, int& i, int upperBound) {
    if (i >= (int)preorder.size() || preorder[i] > upperBound) {
        return nullptr;
    }
    TreeNode* node = new TreeNode(preorder[i++]);
    node->left = build(preorder, i, node->val);
    node->right = build(preorder, i, upperBound);
    return node;
}

TreeNode* bstFromPreorder(vector<int>& preorder) {
    int i = 0;
    return build(preorder, i, INT_MAX);
}

int main() {
    vector<int> preorder = {8, 5, 1, 7, 10, 12};
    TreeNode* root = bstFromPreorder(preorder);
    cout << "Root: " << root->val << endl;
    return 0;
}
\end{lstlisting}

\begin{complexitybox}
\textbf{Time Complexity.} Each element of the preorder array is consumed
exactly once, with $O(1)$ work per element, giving
\[
  O(n).
\]
\textbf{Space Complexity.} $O(h)$ for the recursion stack, plus $O(n)$
for the $n$ nodes created.
\end{complexitybox}

\begin{takeawaybox}
Reconstructing a tree from a single traversal is only possible when that
traversal, combined with some known structural invariant (here, the BST
property), removes all ambiguity about subtree boundaries. The ``pass a
bound downward, consume from a shared pointer'' technique used here is a
direct cousin of Validate BST's (Section~\ref{sec:bst-validate}) bound-passing
recursion, repurposed for construction instead of verification.
\end{takeawaybox}

\section{Recover BST (Two Nodes Swapped)}
\label{sec:bst-recover}

\problemstatement{The values of exactly two nodes in a BST were swapped
by mistake, breaking the BST invariant. Without changing the tree's
structure, recover the original BST by fixing the values of those two
nodes (find them and swap their values back).}

\begin{patternbox}
\emph{``Exactly two nodes are out of place''} in a problem framed around
\textbf{sortedness} is the strongest possible signal to invoke
Section~\ref{sec:bst-intro}'s headline fact directly: a correct BST's
inorder traversal is strictly increasing, so the swapped pair must show
up as exactly one or two places where the inorder sequence decreases
instead of increases -- find those violation points during a single
inorder traversal, and the swapped nodes are sitting right there.
\end{patternbox}

\subsection*{Understanding the problem carefully}

In a correct BST, scanning the inorder sequence left to right, every
value is greater than the one before it. With two values swapped, this
property breaks at either one place (if the two swapped nodes happen to
be \emph{adjacent} in inorder order) or two places (if they are not
adjacent) -- and at each violation, the \emph{earlier} of the two
compared values is the one that is too large for its position, or the
\emph{later} one is too small.

\begin{ideabox}[Track the First and Second Violation During One Inorder Pass]
Maintain $\textit{prev}$ (the previously visited node) and two
not-yet-found pointers, $\textit{first}$ and $\textit{second}$. During
inorder traversal, at each node: if $\textit{prev}$ exists and
$\textit{prev}.\textit{val} > \textit{node}.\textit{val}$ (a violation):
if $\textit{first}$ has not yet been set, set $\textit{first} \gets
\textit{prev}$; in either case, set $\textit{second} \gets
\textit{node}$ (always update to the most recent violation's second
node). Update $\textit{prev} \gets \textit{node}$ and continue. After
the traversal, swap $\textit{first}.\textit{val}$ and
$\textit{second}.\textit{val}$.
\end{ideabox}

\subsection*{Why exactly these two recorded nodes are the swapped pair}

If the swap is between adjacent inorder values, exactly one violation
occurs, and $\textit{first}$/$\textit{second}$ from that single violation
are exactly the two swapped nodes. If the swap is between non-adjacent
values, exactly \emph{two} violations occur: the first violation's
\textit{prev} is the node that ended up holding too large a value (it
should be smaller, since the swap displaced a large value into an
earlier position), and the second violation's current node is the one
holding too small a value (displaced later than it should be). Always
keeping $\textit{first}$ fixed at the \emph{first} violation encountered,
while letting $\textit{second}$ update to track the \emph{most recent}
violation's later node, correctly identifies both endpoints of the
original swap regardless of which of these two cases occurred --
swapping their values back is enough to restore strict monotonicity,
since no third node is ever out of place.

\subsection*{Algorithm}

\begin{enumerate}
  \item Initialize $\textit{prev}, \textit{first}, \textit{second}
        \gets$ null.
  \item Perform an inorder traversal. At each visited node:
  \begin{enumerate}
    \item If $\textit{prev} \ne$ null and $\textit{prev}.\textit{val} >
          \textit{node}.\textit{val}$:
    \begin{itemize}
      \item If $\textit{first} = $ null: $\textit{first} \gets
            \textit{prev}$.
      \item $\textit{second} \gets \textit{node}$.
    \end{itemize}
    \item $\textit{prev} \gets \textit{node}$.
  \end{enumerate}
  \item Swap $\textit{first}.\textit{val}$ and
        $\textit{second}.\textit{val}$.
\end{enumerate}

\subsection*{Dry run}

\dryrun{A BST whose intended inorder sequence is $1,2,3,4,5,6,7$ has
nodes holding values $3$ and $6$ swapped, giving the actual inorder
sequence $1,2,6,4,5,3,7$.}

\begin{itemize}
  \item Visit $1$: $\textit{prev}=$ null, no check. $\textit{prev}=1$.
  \item Visit $2$: $1>2$? No. $\textit{prev}=2$.
  \item Visit $6$: $2>6$? No. $\textit{prev}=6$.
  \item Visit $4$: $6>4$? Yes -- violation. $\textit{first}=$ null, so
        $\textit{first} \gets$ node holding $6$. $\textit{second} \gets$
        node holding $4$. $\textit{prev}=4$.
  \item Visit $5$: $4>5$? No. $\textit{prev}=5$.
  \item Visit $3$: $5>3$? Yes -- violation. $\textit{first}$ already
        set, unchanged. $\textit{second} \gets$ node holding $3$.
        $\textit{prev}=3$.
  \item Visit $7$: $3>7$? No. $\textit{prev}=7$.
\end{itemize}

$\textit{first}$ is the node currently holding $6$; $\textit{second}$ is
the node currently holding $3$. Swap their values: the node that held $6$
now holds $3$, and the node that held $3$ now holds $6$. The inorder
sequence becomes $1,2,3,4,5,6,7$ -- correctly recovered.

\subsection*{C++ implementation}

\begin{lstlisting}
#include <bits/stdc++.h>
using namespace std;

struct TreeNode {
    int val;
    TreeNode* left;
    TreeNode* right;
    TreeNode(int v) : val(v), left(nullptr), right(nullptr) {}
};

void inorder(TreeNode* node, TreeNode*& prev, TreeNode*& first, TreeNode*& second) {
    if (!node) return;

    inorder(node->left, prev, first, second);

    if (prev && prev->val > node->val) {
        if (!first) first = prev;
        second = node;
    }
    prev = node;

    inorder(node->right, prev, first, second);
}

void recoverTree(TreeNode* root) {
    TreeNode *prev = nullptr, *first = nullptr, *second = nullptr;
    inorder(root, prev, first, second);
    swap(first->val, second->val);
}

int main() {
    // Built so that an inorder traversal visits 1, 2, 6, 4, 5, 3, 7 --
    // i.e. the values 3 and 6 have been swapped relative to a correct BST.
    TreeNode* root = new TreeNode(4);
    root->left = new TreeNode(2);
    root->left->left = new TreeNode(1);
    root->left->right = new TreeNode(6); // swapped value
    root->right = new TreeNode(3);       // swapped value
    root->right->left = new TreeNode(5);
    root->right->right = new TreeNode(7);

    recoverTree(root);

    TreeNode *p = nullptr, *f = nullptr, *s = nullptr;
    inorder(root, p, f, s); // verify: no violations should remain
    cout << (f == nullptr ? "Recovered" : "Still broken") << endl;
    return 0;
}
\end{lstlisting}

\begin{complexitybox}
\textbf{Time Complexity.} $O(n)$ -- a single inorder traversal visiting
every node once.
\textbf{Space Complexity.} $O(h)$ for the recursion stack. (An advanced
variant using Morris traversal achieves $O(1)$ space, at the cost of
temporarily modifying the tree's pointers during the scan -- worth
knowing exists, though the recursive version here is the one to derive
first.)
\end{complexitybox}

\begin{takeawaybox}
``Exactly $k$ elements are out of place in an otherwise sorted sequence''
problems are solved by scanning once and recording violation points,
rather than sorting from scratch. Here, $k=2$ gives exactly the
\textit{first}/\textit{second} bookkeeping shown; the same idea
generalizes (with more bookkeeping) to other small, fixed values of $k$.
\end{takeawaybox}

\section{Largest BST in a Binary Tree}
\label{sec:bst-largest-bst}

\problemstatement{Given the root of a (general, not necessarily search)
binary tree, find the size (number of nodes) of the \textbf{largest
subtree} that is a valid BST.}

\begin{patternbox}
This closing problem of the chapter is a deliberate capstone: it combines
Validate BST's (Section~\ref{sec:bst-validate}) range-checking idea with
a \textbf{bottom-up, tree-DP} style of recursion -- each subtree must
report enough information upward (is it a valid BST? how big? what are
its min and max values?) for its \emph{parent} to decide, in $O(1)$
additional work, whether combining it with its sibling and itself forms
an even larger valid BST. The keyword combination \emph{``largest
substructure satisfying a property''} within a \emph{general} tree (not
already known to be a BST) is the signal for this bottom-up,
information-passing recursion style, much like dynamic programming on
trees.
\end{patternbox}

\subsection*{Understanding the problem carefully}

Section~\ref{sec:bst-validate} validated a BST top-down, passing a valid
range \emph{downward}. Here, we do not know in advance which subtrees are
valid BSTs, so we instead compute, for \emph{every} subtree, bottom-up:
whether it is itself a valid BST, its size if so, and its minimum and
maximum values (needed by its parent to check the BST invariant one
level up). A node's own subtree is a valid BST exactly when both children
report ``yes, I am a valid BST,'' \emph{and} the node's value is strictly
greater than its left child's reported maximum, \emph{and} strictly less
than its right child's reported minimum.

\begin{ideabox}[Bottom-Up: Each Subtree Reports (isBST, size, min, max)]
Recurse postorder (process both children before combining at the current
node). For a null node, return $(\textit{isBST}=\text{true},
\textit{size}=0, \textit{min}=+\infty, \textit{max}=-\infty)$ (a
neutral, vacuously-true base case that combines safely with any actual
node). At a real node: if both children report $\textit{isBST}=\text{true}$,
and $\textit{node.val} > \textit{leftInfo.max}$, and
$\textit{node.val} < \textit{rightInfo.min}$: this subtree is a valid
BST of size $\textit{leftInfo.size} + \textit{rightInfo.size} + 1$; update
a running global answer with this size, and return
$(\text{true}, \textit{size}, \min(\textit{leftInfo.min},
\textit{node.val}), \max(\textit{rightInfo.max}, \textit{node.val}))$.
Otherwise, this subtree is not a valid BST: return
$\textit{isBST}=\text{false}$ (its size/min/max no longer matter to any
ancestor, since no ancestor can validly combine across an already-invalid
subtree).
\end{ideabox}

\subsection*{Why the global answer must be tracked throughout, not just at the root}

The largest valid BST is not necessarily the \emph{whole} tree, nor even
necessarily a subtree rooted near the root -- it could be a valid BST
entirely contained within one branch, while everything above that branch
fails the BST invariant. Because the recursion already computes,
\emph{at every single node}, whether that node's subtree is a valid BST
and how large it is, the correct overall answer is simply the
\textbf{maximum} of this size over every node visited during the entire
traversal -- which is why a running global maximum, updated every time a
valid BST subtree is found (regardless of whether its ancestors turn out
to be valid), is necessary, rather than only inspecting the final result
returned at the root.

\subsection*{Algorithm}

\begin{enumerate}
  \item Initialize a global $\textit{maxBSTSize} \gets 0$.
  \item $\textsc{Solve}(\textit{node})$:
  \begin{enumerate}
    \item If $\textit{node} = $ null: return $(\text{true}, 0, +\infty,
          -\infty)$.
    \item $\textit{L} \gets \textsc{Solve}(\textit{node}.\textit{left})$,
          $\textit{R} \gets
          \textsc{Solve}(\textit{node}.\textit{right})$.
    \item If $\textit{L}.\textit{isBST}$ and $\textit{R}.\textit{isBST}$
          and $\textit{node}.\textit{val} > \textit{L}.\textit{max}$ and
          $\textit{node}.\textit{val} < \textit{R}.\textit{min}$:
    \begin{itemize}
      \item $\textit{size} \gets \textit{L}.\textit{size} +
            \textit{R}.\textit{size} + 1$.
      \item $\textit{maxBSTSize} \gets \max(\textit{maxBSTSize},
            \textit{size})$.
      \item Return $(\text{true}, \textit{size},
            \min(\textit{L}.\textit{min}, \textit{node}.\textit{val}),
            \max(\textit{R}.\textit{max}, \textit{node}.\textit{val}))$.
    \end{itemize}
    \item Else: return $(\text{false}, 0, -, -)$.
  \end{enumerate}
  \item Call $\textsc{Solve}(\textit{root})$; return
        $\textit{maxBSTSize}$.
\end{enumerate}

\subsection*{Dry run}

\dryrun{Take the tree with root $5$, left child $2$ (left child $1$,
right child $3$, a valid BST of size $3$), right child $7$ (left child
$6$, right child $4$ -- note $4$ on the right side violates BST order).}

\begin{itemize}
  \item $\textsc{Solve}(1)$: leaf. Returns $(\text{true},1,1,1)$.
  \item $\textsc{Solve}(3)$: leaf. Returns $(\text{true},1,3,3)$.
  \item $\textsc{Solve}(2)$: $L=(\text{true},1,1,1)$,
        $R=(\text{true},1,3,3)$. Check: $2>1$ and $2<3$: valid. Size
        $=1+1+1=3$. $\textit{maxBSTSize}=3$. Return
        $(\text{true},3,1,3)$.
  \item $\textsc{Solve}(6)$: leaf. Returns $(\text{true},1,6,6)$.
  \item $\textsc{Solve}(4)$: leaf. Returns $(\text{true},1,4,4)$.
  \item $\textsc{Solve}(7)$: $L=(\text{true},1,6,6)$ (left child $6$),
        $R=(\text{true},1,4,4)$ (right child $4$). Check: $7>6$? Yes.
        $7<4$? No -- invalid (right child's minimum $4$ is not greater
        than $7$). Return $(\text{false}, -, -, -)$.
  \item $\textsc{Solve}(5)$: $L=(\text{true},3,1,3)$ (from node $2$),
        $R=(\text{false}, \ldots)$ (from node $7$). $R.\textit{isBST}$
        is false, so node $5$'s subtree is not a valid BST. Return
        $(\text{false}, -, -, -)$.
\end{itemize}

The largest valid BST found during the traversal has size $3$ (the
subtree rooted at node $2$), even though the overall tree (rooted at $5$)
is not a valid BST.

\subsection*{C++ implementation}

\begin{lstlisting}
#include <bits/stdc++.h>
using namespace std;

struct TreeNode {
    int val;
    TreeNode* left;
    TreeNode* right;
    TreeNode(int v) : val(v), left(nullptr), right(nullptr) {}
};

struct Info {
    bool isBST;
    int size;
    int minVal, maxVal;
};

int maxBSTSize = 0;

Info solve(TreeNode* node) {
    if (!node) return {true, 0, INT_MAX, INT_MIN};

    Info L = solve(node->left);
    Info R = solve(node->right);

    if (L.isBST && R.isBST && node->val > L.maxVal && node->val < R.minVal) {
        int size = L.size + R.size + 1;
        maxBSTSize = max(maxBSTSize, size);
        return {true, size, min(L.minVal, node->val), max(R.maxVal, node->val)};
    }
    return {false, 0, 0, 0};
}

int largestBSTSubtree(TreeNode* root) {
    maxBSTSize = 0;
    solve(root);
    return maxBSTSize;
}

int main() {
    TreeNode* root = new TreeNode(5);
    root->left = new TreeNode(2);
    root->left->left = new TreeNode(1);
    root->left->right = new TreeNode(3);
    root->right = new TreeNode(7);
    root->right->left = new TreeNode(6);
    root->right->right = new TreeNode(4);

    cout << "Largest BST size: " << largestBSTSubtree(root) << endl;
    return 0;
}
\end{lstlisting}

\begin{complexitybox}
\textbf{Time Complexity.} Each node is visited exactly once, doing
$O(1)$ work combining its children's reported info, giving
\[
  O(n).
\]
This is a substantial improvement over the naive approach of running
Validate BST (Section~\ref{sec:bst-validate}) independently on every
single subtree, which would cost $O(n^2)$.
\textbf{Space Complexity.} $O(h)$ for the recursion stack.
\end{complexitybox}

\begin{takeawaybox}
Whenever a problem asks for the largest/best substructure satisfying some
recursively-checkable property within a general tree, look for a
bottom-up recursion where each subtree reports exactly the information
its parent needs to make an $O(1)$ decision -- validity, size, and
boundary values, in this case. This ``return a small summary struct
upward, combine in $O(1)$ at each parent'' technique is the tree
analogue of dynamic programming, and it is the most general and reusable
idea in this entire chapter.
\end{takeawaybox}

